% week-01_CS2104_2026-2027.tex
% Adapted from latex-template-for-CS2104.tex

% LaTeX report for CS2104 Week 01

%%%%%%%%%%%%%%%%%%%%

\documentclass[12pt]{article}

\usepackage{a4wide}
\usepackage[utf8]{inputenc}
\usepackage{graphicx}
\usepackage{hyperref}
\usepackage{amsmath}
\usepackage{amssymb}

\usepackage{color}
% \definecolor{ltblue}{rgb}{0,0.4,0.4}
\definecolor{ltblue}{rgb}{0,0.0,0.0}
% \definecolor{dkblue}{rgb}{0,0.1,0.6}
\definecolor{dkblue}{rgb}{0,0.0,0.0}
% \definecolor{dkgreen}{rgb}{0,0.35,0}
\definecolor{dkgreen}{rgb}{0,0.0,0}
% \definecolor{dkviolet}{rgb}{0.3,0,0.5}
\definecolor{dkviolet}{rgb}{0.0,0,0.0}
% \definecolor{dkred}{rgb}{0.5,0,0}
\definecolor{dkred}{rgb}{0.0,0,0}

%%%%%%%%%%%%%%%%%%%%

\begin{document}

\title{
  CS2104 Report for Week 01
  \\
  Week 01 Exercise Report
}

\author{
    \textbf{Sim Yue Yang} (A0308298R, e1398738@u.nus.edu) \\
    \textbf{Reuben Thomas} (A0273600N, reuben\_thomas@u.nus.edu) \\
    \textbf{Eugene Long} (A0308383B, e1398823@u.nus.edu) \\
    \textbf{Theodore Lim} (A0266620H, e1091280@u.nus.edu) \\
    \textbf{Truong Kim Mai} (A0301529L, e1355895@u.nus.edu) \\
    \textbf{Lin Haoyue} (A0309004U, haoyue@u.nus.edu) \\
    \textbf{Matthew Allan} (A0305669U, e1375480@u.nus.edu) \\
    \textbf{Kee Yen Cheng} (A0308304N, e1398744@u.nus.edu) \\
    \textbf{Lim Hur} (A0302216X, e1356582@u.nus.edu) \\
    \textbf{Umaiza} (A0321876X, e1524783@u.nus.edu)
}

\date{\today}

\maketitle

\begin{abstract}
This document contains the solutions to the Week 01 exercise. The code is provided in a separate file along with this document.
\end{abstract}

\vfill

\begin{center}
...a picture, maybe...
\end{center}

\vfill

\begin{flushright}
\emph{...a relevant quote, possibly...} \\
-- the author of the quote
\end{flushright}

\newpage

\tableofcontents

\newpage

%%%%%%%%%%

\section{Overall Introduction}
\label{sec:overall-introduction}

\begin{itemize}

\item
  This document is the document that documents the solutions for the questions in the Week01 handin document.

\item
  This document documents the solutions for the questions in the Week01 handin.

\item
  This point is situated in week 01 of the semester.

\item
  Title, abstract, content table, introduction, exercises (question, introduction, solution, conclusion for each exercise) \\ and an overall conclusion.

\end{itemize}

%%%%%%%%%%

\newpage

\section{Exercise 1: Self-reference}

%%%%%

\subsection{Statement}

The first lecture started with the DSB jingles \#1 and \#2 at
\href{https://toglyde.dk/toglyde/stationslyde-inkl-hoejttalerudkald/}{toglyde.dk}.

Using the
\href{https://en.wikipedia.org/wiki/Solf\%C3\%A8ge\#Fixed_do_solf\%C3\%A8ge}{fixed-do solf\`ege},
this three-note jingle D/Eb/Bb (pronounced D/S/B in Danish, i.e., D Sharp B) is
the acronym of DSB, the Danish Railroad Company ``Danske Statsbaner'', and is
played in Danish rail stations when a train arrives. It is due to the Danish
composer
\href{https://en.wikipedia.org/wiki/Niels_Viggo_Bentzon}{Niels Viggo Bentzon}.

This jingle illustrates self-reference since reading aloud yields the name DSB.

Other examples of self-reference include the
\href{https://en.wikipedia.org/wiki/Liar_paradox}{liar paradox} and Ren\'e
Magritte's painting
\href{https://en.wikipedia.org/wiki/Ceci_n\%27est_pas_une_pipe}{``The Treachery of Images''}.

\textbf{1.a.} Identify the self-reference in the cover of blues guitarist Coco
Montoya's album
\href{https://cocomontoya.bandcamp.com/album/writing-on-the-wall}{``The Writing on the Wall''}.

\textbf{1.b.} Identify the self-reference in the cover of Gail Carson Levine's
book
\href{https://www.goodreads.com/book/show/136218.Writing_Magic}{``Writing Magic: Creating Stories that Fly''}.

\textbf{1.c.} Identify the self-reference in
\href{http://en.wikipedia.org/wiki/Alan_Perlis}{Alan Perlis}'s
\href{http://www.cs.yale.edu/homes/perlis-alan/quotes.html}{programming epigram \#59}:

\begin{quote}
In English, every word can be verbed.
\end{quote}

All the self-references above are direct. Some self-references are indirect in
that something named A refers to something named B that refers to A. For
example, for any natural number n, the number n+1 is even if the number n is
odd, and the number n+1 is odd if the number n is even.

\textbf{1.d.} Identify the indirect self-reference in the following story:

\begin{quote}
Two women walk into a bar and talk about the
\href{https://en.wikipedia.org/wiki/Bechdel_test}{Bechdel-Wallace test}.
\end{quote}

\textbf{1.e.} Identify the self-reference, if there is one, in the following
sentences:

\begin{itemize}
\item \textbf{1.e.1.} This sentence no verb.
\item \textbf{1.e.2.} Nostalgia isn't what it used to be.
\item \textbf{1.e.3.} Fashion is what goes out of fashion. ---
  \href{https://en.wikipedia.org/wiki/Coco_Chanel}{Coco Chanel}
\item \textbf{1.e.4.} To be or not to be.
  (\href{https://en.wikipedia.org/wiki/To_be_or_not_to_be}{phrase} $\cdot$
  \href{https://en.wikipedia.org/wiki/To_Be_or_Not_to_Be_(1942_film)}{1942 film})
\end{itemize}

%%%%%

\subsection{Introduction}

\begin{itemize}

\item
  The point of the exercise is to identify self-references across various media and sentences.

\item
  The exercise is solved by inspecting each example and describing how it refers to itself.

\end{itemize}

%%%%%

\subsection{Solution}

\textbf{1.a.} The album cover features a literal writing on the wall.

\textbf{1.b.} The person on the cover is writing the word ``Magic'', taking the
phrase ``Writing Magic'' literally.

\textbf{1.c.} The word ``verb'' is used as a verb in the sentence.

\textbf{1.d.} The Bechdel-Wallace test checks if 2 female characters in a story
converse about something other than a male character. This story about two women
conversing about the Bechdel-Wallace test itself passes the Bechdel-Wallace
test.

\textbf{1.e.1} None.

\textbf{1.e.2} The subject is nostalgic for the concept of nostalgia.

\textbf{1.e.3} Fashion is ever-changing, so all fashion can eventually become
irrelevant (``goes out of fashion'').

\textbf{1.e.4} None.

%%%%%

\subsection{Conclusion}

\begin{itemize}

\item
  The point of the exercise was to practice spotting self-reference.

\item
  The learning points illustrated were direct self-reference, and indirect self-reference where A refers to B which refers back to A.

\end{itemize}

%%%%%%%%%%

\newpage

\section{Exercise 2: Media involving media}

%%%%%

\subsection{Statement}

Can you think of

\begin{itemize}
\item a book involving a book,
\item a theater play involving a theater play, or
\item a movie involving a movie?
\end{itemize}

More generally, can you think of one of these mediums (book, theater play,
movie, TV series, etc.) involving one of these mediums?

If so, briefly describe a couple of them and their involvement.

\textbf{Subsidiary question:} Give an example where self-reference is used in
Computer Science.

\textbf{Partial solution for Exercise 2.}
\href{https://en.wikipedia.org/wiki/A_Midsummer_Night\%27s_Dream}{``A Midsummer Night's Dream''}
qualifies because in Scene 1 of Act V,
\href{https://en.wikipedia.org/wiki/Pyramus_and_Thisbe\#Adaptations}{Pyramus and Thisbe}
is a play within a play.

%%%%%

\subsection{Introduction}

\begin{itemize}

\item
  The point of the exercise is to find a medium that contains another instance of a medium, and an example of self-reference in Computer Science.

\item
  The exercise is solved by recalling examples from memory.

\end{itemize}

%%%%%

\subsection{Solution}

\begin{itemize}
\item In Ridge Racer 7, the playable mini-game that appears during loading
  screens is Galaxian, the classic Namco arcade shooter.
\item In computer science, recursive functions are functions that call
  themselves as part of their execution.
\end{itemize}

%%%%%

\subsection{Conclusion}

\begin{itemize}

\item
  The point of the exercise was to spot self-reference in nested media and Computer Science.

\item
  The exercise was solved by giving real-life examples.

\item
  The learning point illustrated was that recursion is a form of self-reference.

\end{itemize}

%%%%%%%%%%

\newpage

\section{Exercise 3: A multiple-choice question}

%%%%%

\subsection{Statement}

Facetiously, you are asked whether there exists a word in English that contains
all of the letters ``a'', ``e'', ``i'', ``o'', ``u'', and ``y'', in that order.
Does there exist such a word in English?

Possible answers:

\begin{itemize}
\item \textbf{a.} Yes, unquestionably. Constant vigilance!
\item \textbf{b.} No.
\item \textbf{c.} I cannot possibly know this, and furthermore this facetious
  question is nefariously out of scope in CS2104.
\end{itemize}

%%%%%

\subsection{Introduction}

\begin{itemize}

\item
  The point of the exercise is to determine whether an English word contains the letters a, e, i, o, u, y in that order.

\item
  The exercise is solved by finding one such word.

\end{itemize}

%%%%%

\subsection{Solution}

\begin{itemize}
\item a. ``Adventitiously''.
\end{itemize}

%%%%%

\subsection{Conclusion}

\begin{itemize}

\item
  The point of the exercise was to ask Google which words contain all vowels in the specified order.

\item
  The exercise was solved by finding the word "adventitiously".

\item
  The learning point illustrated was that the word "adventitiously" exists.

\end{itemize}

%%%%%%%%%%

\newpage

\section{Exercise 4: The Sorcerer's Apprentice}

%%%%%

\subsection{Statement}

In reference to LLMs, please view the short piece ``The Sorcerer's Apprentice''
in \href{https://en.wikipedia.org/wiki/Fantasia_(1940_film)}{``Fantasia''}.

\textbf{N.B.:} While you are at it, expand your musical culture with Paul
Dukas's score, as a continuation of Johann Strauss II's waltz that you so
magnificently rendered at the beginning of the first lecture
(\href{https://en.wikipedia.org/wiki/The_Blue_Danube}{The Blue Danube}).

%%%%%

\subsection{Introduction}

\begin{itemize}

\item
  The point of the exercise is to relate "The Sorcerer's Apprentice" to LLMs.

\item
  The exercise is solved by reading the short piece.

\end{itemize}

%%%%%

\subsection{Solution}

\begin{itemize}
\item The story is about the apprentice using magic, losing control of it, which results in a mess. If we view this as analogy where the apprentice is a CS undergrad and the magic is the LLM, we could say that a CS undergrad using LLMs wantonly will get nowhere. 
\end{itemize}

%%%%%

\subsection{Conclusion}

\begin{itemize}

\item
  The point of the exercise was to draw an analogy between LLMs and magic.

\item
  The exercise was solved by reading the piece.

\item
  The learning point illustrated was the risk of using LLMs without knowledge over its outputs.

\end{itemize}

%%%%%%%%%%

\newpage

\section{Exercise 5: Implementing the Sigma notation using a for loop}

%%%%%

\subsection{Statement}

See \href{https://en.wikipedia.org/wiki/Capital-sigma_notation}{Capital-sigma notation}.

The Sigma notation reads, given a non-negative integer n and a function f from
(non-negative) integers to integers,

\begin{verbatim}
sum for i = 0 to n of f(i)
\end{verbatim}

and it stands for

\begin{verbatim}
f(0) + f(1) + ... + f(n)
\end{verbatim}

In this notation, 0 is the lower bound, n is the upper bound, and f(i) is the
body.

Here, the notation is extended to negative integers in that given a negative
integer m and a function f,

\begin{verbatim}
sum for i = 0 to m of f(i)
\end{verbatim}

stands for

\begin{verbatim}
0
\end{verbatim}

The starting point of this exercise is to implement --- in the programming
language of your choice --- a routine / function / procedure / method /
what-have-you named \texttt{Sigma0} that takes two parameters: an integer n and
a function from integer to integer f. Applying \texttt{Sigma0} to n and f should
yield the result of \texttt{f(0) + f(1) + ... + f(n)} when \texttt{n} is
non-negative, and \texttt{0} otherwise.

\textbf{5.a.} Devise meaningful tests for your implementation. For example,
\texttt{f} could be a constant function, or the identity function, or the
function that maps an integer \texttt{i} to the \texttt{i}+first odd natural
number (namely \texttt{i.2*i + 1} using the lambda notation; then what is
computed is \texttt{1 + 3 + 5 + ...}).

\textbf{5.b.} Implement \texttt{Sigma0} using a for loop.

\textbf{5.c.} Check that your implementation passes your unit tests. (The testing
should be automatic, and since what goes well should go without saying, it
should only report the tests that failed.)

\textbf{5.d.} Implement a function named \texttt{Sigma1} that takes two
parameters: an integer n and a function from integer to integer \texttt{f}.
Applying \texttt{Sigma1} to \texttt{n} and \texttt{f} should yield the result of
\texttt{f(1) + ... + f(n)} when \texttt{n} is positive, and \texttt{0} otherwise.

%%%%%

\subsection{Introduction}

\begin{itemize}

\item
  The point of the exercise is to implement the Sigma notation with a for loop and test it.

\item
  The exercise is solved by accumulating $f(i)$ over the range in a loop, returning 0 for a negative upper bound.

\end{itemize}

%%%%%

\subsection{Solution}

See functions Sigma0 and Sigma1 in week01.cpp.

%%%%%

\subsection{Conclusion}

\begin{itemize}

\item
  The point of the exercise was to implement Sigma notation as code and verify it.

\item
  The exercise was solved with a for-loop accumulator guarding negative bounds, checked by automated tests.

\item
  The learning points illustrated were summation as iteration, higher-order functions, and automated unit testing.

\end{itemize}

%%%%%%%%%%

\newpage

\section{Exercise 6: Sums of floor divisions}

%%%%%

\subsection{Statement}

\textbf{6.a.} Given a non-negative integer n, what is the result of

\begin{verbatim}
sum for i = 0 to 2 * n of floor(n/2)
\end{verbatim}

where floor(5/2) = 2 and floor(6/2) = 3.

\textbf{6.b.} Given a non-negative integer n, what is the result of

\begin{verbatim}
sum for i = 0 to 3 * n of floor(n/3)
\end{verbatim}

\textbf{6.c.} Given a non-negative integer n, what is the result of

\begin{verbatim}
sum for i = 0 to 4 * n of floor(n/4)
\end{verbatim}

Hint: Use the
\href{https://oeis.org/}{Online Encyclopedia of Integer Sequences}.

\textbf{6.d.} Given a non-negative integer n, what is the result of

\begin{verbatim}
sum for i = 0 to 1 * n of floor(n/1)
\end{verbatim}

(Can you simplify this formulation?)

%%%%%

\subsection{Introduction}

\begin{itemize}

\item
  The point of the exercise is to find closed forms for sums whose body is a constant floor term.

\item
  The exercise is solved by noting the body does not depend on $i$, so the result is the body times the number of terms.

\end{itemize}

%%%%%

\subsection{Solution}

For each part let $r$ be the result.

\textbf{6.a.} $r = (2n + 1)\left\lfloor \dfrac{n}{2} \right\rfloor$

\textbf{6.b.} $r = (3n + 1)\left\lfloor \dfrac{n}{3} \right\rfloor$

\textbf{6.c.} $r = (4n + 1)\left\lfloor \dfrac{n}{4} \right\rfloor$

\textbf{6.d.} $r = (n + 1)(n)$

%%%%%

\subsection{Conclusion}

\begin{itemize}

\item
  The point of the exercise was to derive closed forms for constant-body sums.

\item
  The exercise was solved by multiplying the number of terms by the constant body.

\item
  The learning point illustrated was that a constant summand gives (count) $\times$ (value).

\end{itemize}

%%%%%%%%%%

\newpage

\section{Exercise 7: Nested sums and binomials}

%%%%%

\subsection{Statement}

\textbf{7.a.} What is the result of

\begin{verbatim}
1
\end{verbatim}

\textbf{7.b.} Given a non-negative integer x, what is the result of

\begin{verbatim}
sum for i1 = 0 to x of
1
\end{verbatim}

\textbf{7.c.} Given a non-negative integer x, what is the result of

\begin{verbatim}
sum for i2 = 0 to x of
sum for i1 = 0 to i2 of
1
\end{verbatim}

\textbf{7.d.} Given a non-negative integer x, what is the result of

\begin{verbatim}
sum for i3 = 0 to x of
sum for i2 = 0 to i3 of
sum for i1 = 0 to i2 of
1
\end{verbatim}

\textbf{7.e.} Given a non-negative integer x, what is the result of

\begin{verbatim}
sum for i4 = 0 to x of
sum for i3 = 0 to i4 of
sum for i2 = 0 to i3 of
sum for i1 = 0 to i2 of
1
\end{verbatim}

\textbf{7.f.} Implement a function that maps two non-negative integers N and x to

\begin{verbatim}
sum for iN = 0 to x of
sum for iN-1 = 0 to iN of
...
sum for i2 = 0 to i3 of
sum for i1 = 0 to i2 of
1
\end{verbatim}

$\displaystyle\sum_{i_{N}=0}^{x}\sum_{i_{N-1}=0}^{i_{N}}\dots\sum_{i_{2}=0}^{i_{3}}\sum_{i_{1}=0}^{i_{2}}1$

and empirically characterize its result. If possible (ha, ha, ha) \\ justify your
characterization.

\textbf{N.B.:} Since there is no ambiguity, feel free to write

\begin{verbatim}
sum for i = 0 to x of
sum for i = 0 to i of
...
sum for i = 0 to i of
sum for i = 0 to i of
1
\end{verbatim}

but then you need to say that there are N occurrences of ``sum'' since this
information is no longer present in the notation.

%%%%%

\subsection{Introduction}

\begin{itemize}

\item
  The point of the exercise is to characterize the result of $N$ nested sums of 1.

\item
  The exercise is solved by computing values into a sequence, and looking them up on OEIS.

\end{itemize}

%%%%%

\subsection{Solution}

For each part let $r$ be the result.

\textbf{7.a.} $r = 1$

\textbf{7.b.} $r = x + 1$

\textbf{7.c.} $r = \dfrac{(x + 1)(x + 2)}{2}$

\textbf{7.d.} $r = \dfrac{(x + 1)(x + 2)(x + 3)}{6}$

\textbf{7.e.} $r = \dfrac{(x + 1)(x + 2)(x + 3)(x + 4)}{24}$

\textbf{7.f.}

See function q7 in week01.cpp.

The following are the outputs for the function for $N \in [0,5]$ and
$x \in [0,5]$:

\begin{verbatim}
N = 0, x = 0 | 1
N = 0, x = 1 | 1
N = 0, x = 2 | 1
N = 0, x = 3 | 1
N = 0, x = 4 | 1
N = 0, x = 5 | 1
N = 1, x = 0 | 1
N = 1, x = 1 | 2
N = 1, x = 2 | 3
N = 1, x = 3 | 4
N = 1, x = 4 | 5
N = 1, x = 5 | 6
N = 2, x = 0 | 1
N = 2, x = 1 | 3
N = 2, x = 2 | 6
N = 2, x = 3 | 10
N = 2, x = 4 | 15
N = 2, x = 5 | 21
N = 3, x = 0 | 1
N = 3, x = 1 | 4
N = 3, x = 2 | 10
N = 3, x = 3 | 20
N = 3, x = 4 | 35
N = 3, x = 5 | 56
N = 4, x = 0 | 1
N = 4, x = 1 | 5
N = 4, x = 2 | 15
N = 4, x = 3 | 35
N = 4, x = 4 | 70
N = 4, x = 5 | 126
N = 5, x = 0 | 1
N = 5, x = 1 | 6
N = 5, x = 2 | 21
N = 5, x = 3 | 56
N = 5, x = 4 | 126
N = 5, x = 5 | 252
\end{verbatim}

I have manually calculated the results and verified them to be accurate.

By searching the OEIS, these seem to be the binomial results of
$\displaystyle\binom{x + N}{N}$.

It can be justified by observing that
$\displaystyle\binom{a}{b} = \dfrac{a \cdot (a - 1) \cdot \dots \cdot (a - b + 1)}{k \cdot (k - 1) \cdot \dots \cdot 1}$,
so all that's needed is to offset the $x$ by $+N$ in the top part.

%%%%%

\subsection{Conclusion}

\begin{itemize}

\item
  The point of the exercise was to identify the value of $N$ nested sums of 1.

\item
  The exercise was solved by recognizing outputs as $\binom{x + N}{N}$.

\item
  The learning point illustrated was that nested sums count combinations.

\end{itemize}

%%%%%%%%%%

\newpage

\section{Exercise 8: Catalan numbers}

%%%%%

\subsection{Statement}

Implement a function that maps non-negative integer N to

\begin{verbatim}
sum for iN = 0 to 0 of
sum for iN-1 = 0 to iN+1 of
...
sum for i2 = 0 to i3+1 of
sum for i1 = 0 to i2+1 of
1
\end{verbatim}

$\displaystyle\sum_{i_{N}=0}^{0}\sum_{i_{N-1}=0}^{i_{N}+1}\dots\sum_{i_{2}=0}^{i_{3}+1}\sum_{i_{1}=0}^{i_{2}+1}1$

and empirically characterize its result. If possible (ha, ha, ha) \\ justify your
characterization.

%%%%%

\subsection{Introduction}

\begin{itemize}

\item
  The point of the exercise is to characterize a nested sum whose bounds are shifted by $+1$.

\item
  The exercise is solved by a recursive implementation, tabulating the outputs, and matching against the OEIS.

\end{itemize}

%%%%%

\subsection{Solution}

See function q8 in week01.cpp.

The following are the outputs for the function for $N \in [0,10]$:

\begin{verbatim}
N = 0 | 1
N = 1 | 1
N = 2 | 2
N = 3 | 5
N = 4 | 14
N = 5 | 42
N = 6 | 132
N = 7 | 429
N = 8 | 1430
N = 9 | 4862
N = 10 | 16796
\end{verbatim}

I have manually calculated the results and verified them to be accurate.

By searching the OEIS, these seem to be the Catalan sequence.

After reading the Wikipedia page on Catalan numbers, I have concluded that I do
not understand it, and therefore will not attempt to justify the
characterization.

%%%%%

\subsection{Conclusion}

\begin{itemize}

\item
  The point of the exercise was to identify the value of the shifted nested sum.

\item
  The exercise was solved by recursion and an OEIS lookup, giving the Catalan numbers.

\item
  The learning point illustrated was the Catalan sequence.

\end{itemize}

%%%%%%%%%%

\newpage

\section{Exercise 9: Fibonacci sequence}

%%%%%

\subsection{Statement}

Implement a function that maps non-negative integer N

\begin{verbatim}
sum for iN = 0 to 0 of
sum for iN-1 = 0 to 1-iN of
...
sum for i2 = 0 to 1-i3 of
sum for i1 = 0 to 1-i2 of
1
\end{verbatim}

$\displaystyle\sum_{i_{N}=0}^{0}\sum_{i_{N-1}=0}^{1-i_{N}}\dots\sum_{i_{2}=0}^{1-i_{3}}\sum_{i_{1}=0}^{1-i_{2}}1$

and characterize its result. (Remember that if the upper bound of a summation is
negative, the result is 0.) If possible (ha, ha, ha) \\ justify your
characterization.

%%%%%

\subsection{Introduction}

\begin{itemize}

\item
  The point of the exercise is to characterize a nested sum whose bounds are of the form $1 - i$.

\item
  The exercise is solved by a recursive implementation, getting the outputs, and matching against the OEIS.

\end{itemize}

%%%%%

\subsection{Solution}

See function q9 in week01.cpp.

The following are the outputs for the function for $N \in [0,10]$:

\begin{verbatim}
N = 0 | 1
N = 1 | 1
N = 2 | 2
N = 3 | 3
N = 4 | 5
N = 5 | 8
N = 6 | 13
N = 7 | 21
N = 8 | 34
N = 9 | 55
N = 10 | 89
\end{verbatim}

I have manually calculated the results and verified them to be accurate.

By searching the OEIS, these seem to be the $N$-th number of the fibonacci
sequence.

It can be justified by observing that to calculate for $N$, the count of the 2nd
sigma operation is between 0 and 1. When that count is 0, it recursively
calculates for $N-1$, and when the count is 1, it recursively calculates for
$N-2$. Then the 2nd sigma operation simply sums up the result.

%%%%%

\subsection{Conclusion}

\begin{itemize}

\item
  The point of the exercise was to identify the value of this nested sum.

\item
  The exercise was solved by recursion producing $f(N-1) + f(N-2)$, giving the Fibonacci numbers.

\item
  The learning point illustrated was the Fibonacci recurrence arising from the summation bounds.

\end{itemize}

%%%%%%%%%%

\newpage

\section{Exercise 10: Factorials (decreasing bounds)}

%%%%%

\subsection{Statement}

Implement a function that maps non-negative integer N

\begin{verbatim}
sum for i = 0 to N of
sum for i = 0 to N-1 of
...
sum for i = 0 to 2 of
sum for i = 0 to 1 of
sum for i = 0 to 0 of
1
\end{verbatim}

$\displaystyle\sum_{i=0}^{N}\sum_{i=0}^{N-1}\dots\sum_{i=0}^{2}\sum_{i=0}^{1}\sum_{i=0}^{0}1$

and characterize its result. Justify your characterization.

%%%%%

\subsection{Introduction}

\begin{itemize}

\item
  The point of the exercise is to characterize nested sums with constant bounds decreasing from $N$ to 0.

\item
  The exercise is solved by multiplying the layer sizes, getting the outputs, and matching against the OEIS.

\end{itemize}

%%%%%

\subsection{Solution}

See function q10 in week01.cpp.

The following are the outputs for the function for $N \in [0,10]$:

\begin{verbatim}
N = 0 | 1
N = 1 | 2
N = 2 | 6
N = 3 | 24
N = 4 | 120
N = 5 | 720
N = 6 | 5040
N = 7 | 40320
N = 8 | 362880
N = 9 | 3628800
N = 10 | 39916800
\end{verbatim}

I have manually calculated the results and verified them to be accurate.

By searching the OEIS, these seem to be the $N + 1$ factorials.

It can be justified by observing that each Sigma operation nests another Sigma
operation within that decreases the upper bound by 1, from $N$ to 0. Since it's
all just adding up 1s, the final sum can be calculated by finding the number of
total rounds by all Sigma operations. And the multiplied terms are just the
factorial operation $(N + 1)!$ in reverse.

%%%%%

\subsection{Conclusion}

\begin{itemize}

\item
  The point of the exercise was to identify the value of these nested sums.

\item
  The exercise was solved by taking the product of the layer sizes, giving $(N + 1)!$.

\item
  The learning point illustrated was factorial as a product of counts.

\end{itemize}

%%%%%%%%%%

\newpage

\section{Exercise 11: Factorials (increasing bounds)}

%%%%%

\subsection{Statement}

Implement a function that maps non-negative integer N

\begin{verbatim}
sum for i = 0 to 0 of
sum for i = 0 to 1 of
...
sum for i = 0 to N-1 of
sum for i = 0 to N of
1
\end{verbatim}

$\displaystyle\sum_{i=0}^{0}\sum_{i=0}^{1}\dots\sum_{i=0}^{N-1}\sum_{i=0}^{N}1$

and characterize its result. Justify your characterization.

%%%%%

\subsection{Introduction}

\begin{itemize}

\item
  The point of the exercise is to characterize nested sums with constant bounds increasing from 0 to $N$.

\item
  The exercise is solved by multiplying the layer sizes, getting the outputs, and matching against the OEIS.

\end{itemize}

%%%%%

\subsection{Solution}

See function q11 in week01.cpp.

The following are the outputs for the function for $N \in [0,10]$:

\begin{verbatim}
N = 0 | 1
N = 1 | 2
N = 2 | 6
N = 3 | 24
N = 4 | 120
N = 5 | 720
N = 6 | 5040
N = 7 | 40320
N = 8 | 362880
N = 9 | 3628800
N = 10 | 39916800
\end{verbatim}

I have manually calculated the results and verified them to be accurate.

By searching the OEIS, these seem to be the $N + 1$ factorials.

It can be justified by observing that each Sigma operation nests another Sigma
operation within that increases the upper bound by 1, from 0 up to $N$. Since
it's all just adding up 1s, the final sum can be calculated by finding the
number of total rounds by all Sigma operations. And the multiplied terms are
just the factorial operation $(N + 1)!$ in reverse.

%%%%%

\subsection{Conclusion}

\begin{itemize}

\item
  The point of the exercise was to identify the value of these nested sums.

\item
  The exercise was solved by taking the product of the layer sizes, giving $(N + 1)!$ as in Exercise 10.

\item
  The learning point illustrated was that the order of the bounds does not change the product.

\end{itemize}

%%%%%%%%%%

\newpage

\section{Food for Thought}

%%%%%

\subsection{Statement}

These nested sums are not efficient. (Why?) How could their efficiency be
improved?

%%%%%

\subsection{Solution}

These nested sums, for a certain $N > 0$ requires the $N - 1$ levels of the
nested sum to be computed before it can continue. The efficiency can be improved
by finding a formula that allows for direct computation.

%%%%%%%%%%

\newpage

\section{Overall Conclusion}

The overall conclusion goes here:

\begin{itemize}

\item
  We are in week 01 of the semester and the general topic addressed in this report is self-reference in both media and in mathematics / computer science (specifically Sigma / for loops).

\item
  Identification of self-referencing, as well as recognition of its results.

\end{itemize}

%%%%%%%%%%%%%%%%%%%%

\end{document}
